% figs.tex —— 由 ch11/make_figs.py 產生，請勿手動修改。
% 重新產生：python ch11/make_figs.py（第 11 章，動畫 W1101–W1106）

\newcommand{\ElevenFigKoch}{%
\begin{tikzpicture}[font=\scriptsize,baseline=(current bounding box.center)]
\draw[cBar,line join=round] (0.000,0.000) -- (2.950,0.000);
\node[anchor=north] at (1.475,-0.1) {$k=0$};
\draw[cBar,line join=round] (3.280,0.000) -- (4.263,0.000) -- (4.755,0.852) -- (5.247,0.000) -- (6.230,0.000);
\node[anchor=north] at (4.755,-0.1) {$k=1$};
\draw[cBar,line join=round] (6.560,0.000) -- (6.888,0.000) -- (7.052,0.284) -- (7.216,0.000) -- (7.543,0.000) -- (7.707,0.284) -- (7.543,0.568) -- (7.871,0.568) -- (8.035,0.852) -- (8.199,0.568) -- (8.527,0.568) -- (8.363,0.284) -- (8.527,0.000) -- (8.854,0.000) -- (9.018,0.284) -- (9.182,0.000) -- (9.510,0.000);
\node[anchor=north] at (8.035,-0.1) {$k=2$};
\draw[cBar,line join=round] (9.840,0.000) -- (9.949,0.000) -- (10.004,0.095) -- (10.059,0.000) -- (10.168,0.000) -- (10.222,0.095) -- (10.168,0.189) -- (10.277,0.189) -- (10.332,0.284) -- (10.386,0.189) -- (10.496,0.189) -- (10.441,0.095) -- (10.496,0.000) -- (10.605,0.000) -- (10.659,0.095) -- (10.714,0.000) -- (10.823,0.000) -- (10.878,0.095) -- (10.823,0.189) -- (10.933,0.189) -- (10.987,0.284) -- (10.933,0.378) -- (10.823,0.378) -- (10.878,0.473) -- (10.823,0.568) -- (10.933,0.568) -- (10.987,0.662) -- (11.042,0.568) -- (11.151,0.568) -- (11.206,0.662) -- (11.151,0.757) -- (11.260,0.757) -- (11.315,0.852) -- (11.370,0.757) -- (11.479,0.757) -- (11.424,0.662) -- (11.479,0.568) -- (11.588,0.568) -- (11.643,0.662) -- (11.697,0.568) -- (11.807,0.568) -- (11.752,0.473) -- (11.807,0.378) -- (11.697,0.378) -- (11.643,0.284) -- (11.697,0.189) -- (11.807,0.189) -- (11.752,0.095) -- (11.807,0.000) -- (11.916,0.000) -- (11.971,0.095) -- (12.025,0.000) -- (12.134,0.000) -- (12.189,0.095) -- (12.134,0.189) -- (12.244,0.189) -- (12.298,0.284) -- (12.353,0.189) -- (12.462,0.189) -- (12.408,0.095) -- (12.462,0.000) -- (12.571,0.000) -- (12.626,0.095) -- (12.681,0.000) -- (12.790,0.000);
\node[anchor=north] at (11.315,-0.1) {$k=3$};
\begin{scope}
\clip (13.070,-0.12) rectangle (16.190,0.980);
\draw[cExact,densely dashed] (13.120,0) circle (0.983);
\draw[gray,densely dotted] (13.120,0) circle (2.950);
\end{scope}
\draw[cBar,line join=round] (14.103,0.000) -- (14.122,0.032) -- (14.103,0.063) -- (14.140,0.063) -- (14.158,0.095) -- (14.140,0.126) -- (14.103,0.126) -- (14.122,0.158) -- (14.103,0.189) -- (14.140,0.189) -- (14.158,0.221) -- (14.176,0.189) -- (14.213,0.189) -- (14.231,0.221) -- (14.213,0.252) -- (14.249,0.252) -- (14.267,0.284) -- (14.249,0.315) -- (14.213,0.315) -- (14.231,0.347) -- (14.213,0.378) -- (14.176,0.378) -- (14.158,0.347) -- (14.140,0.378) -- (14.103,0.378) -- (14.122,0.410) -- (14.103,0.442) -- (14.140,0.442) -- (14.158,0.473) -- (14.140,0.505) -- (14.103,0.505) -- (14.122,0.536) -- (14.103,0.568) -- (14.140,0.568) -- (14.158,0.599) -- (14.176,0.568) -- (14.213,0.568) -- (14.231,0.599) -- (14.213,0.631) -- (14.249,0.631) -- (14.267,0.662) -- (14.285,0.631) -- (14.322,0.631) -- (14.304,0.599) -- (14.322,0.568) -- (14.358,0.568) -- (14.376,0.599) -- (14.395,0.568) -- (14.431,0.568) -- (14.449,0.599) -- (14.431,0.631) -- (14.468,0.631) -- (14.486,0.662) -- (14.468,0.694) -- (14.431,0.694) -- (14.449,0.725) -- (14.431,0.757) -- (14.468,0.757) -- (14.486,0.789) -- (14.504,0.757) -- (14.540,0.757) -- (14.559,0.789) -- (14.540,0.820) -- (14.577,0.820) -- (14.595,0.852) -- (14.613,0.820) -- (14.650,0.820) -- (14.631,0.789) -- (14.650,0.757) -- (14.686,0.757) -- (14.704,0.789) -- (14.722,0.757) -- (14.759,0.757) -- (14.741,0.725) -- (14.759,0.694) -- (14.722,0.694) -- (14.704,0.662) -- (14.722,0.631) -- (14.759,0.631) -- (14.741,0.599) -- (14.759,0.568) -- (14.795,0.568) -- (14.814,0.599) -- (14.832,0.568) -- (14.868,0.568) -- (14.886,0.599) -- (14.868,0.631) -- (14.905,0.631) -- (14.923,0.662) -- (14.941,0.631) -- (14.977,0.631) -- (14.959,0.599) -- (14.977,0.568) -- (15.014,0.568) -- (15.032,0.599) -- (15.050,0.568) -- (15.087,0.568) -- (15.068,0.536) -- (15.087,0.505) -- (15.050,0.505) -- (15.032,0.473) -- (15.050,0.442) -- (15.087,0.442) -- (15.068,0.410) -- (15.087,0.378) -- (15.050,0.378) -- (15.032,0.347) -- (15.014,0.378) -- (14.977,0.378) -- (14.959,0.347) -- (14.977,0.315) -- (14.941,0.315) -- (14.923,0.284) -- (14.941,0.252) -- (14.977,0.252) -- (14.959,0.221) -- (14.977,0.189) -- (15.014,0.189) -- (15.032,0.221) -- (15.050,0.189) -- (15.087,0.189) -- (15.068,0.158) -- (15.087,0.126) -- (15.050,0.126) -- (15.032,0.095) -- (15.050,0.063) -- (15.087,0.063) -- (15.068,0.032) -- (15.087,0.000) -- (15.123,0.000) -- (15.141,0.032) -- (15.160,0.000) -- (15.196,0.000) -- (15.214,0.032) -- (15.196,0.063) -- (15.232,0.063) -- (15.251,0.095) -- (15.269,0.063) -- (15.305,0.063) -- (15.287,0.032) -- (15.305,0.000) -- (15.342,0.000) -- (15.360,0.032) -- (15.378,0.000) -- (15.414,0.000) -- (15.433,0.032) -- (15.414,0.063) -- (15.451,0.063) -- (15.469,0.095) -- (15.451,0.126) -- (15.414,0.126) -- (15.433,0.158) -- (15.414,0.189) -- (15.451,0.189) -- (15.469,0.221) -- (15.487,0.189) -- (15.524,0.189) -- (15.542,0.221) -- (15.524,0.252) -- (15.560,0.252) -- (15.578,0.284) -- (15.597,0.252) -- (15.633,0.252) -- (15.615,0.221) -- (15.633,0.189) -- (15.669,0.189) -- (15.688,0.221) -- (15.706,0.189) -- (15.742,0.189) -- (15.724,0.158) -- (15.742,0.126) -- (15.706,0.126) -- (15.688,0.095) -- (15.706,0.063) -- (15.742,0.063) -- (15.724,0.032) -- (15.742,0.000) -- (15.779,0.000) -- (15.797,0.032) -- (15.815,0.000) -- (15.851,0.000) -- (15.870,0.032) -- (15.851,0.063) -- (15.888,0.063) -- (15.906,0.095) -- (15.924,0.063) -- (15.961,0.063) -- (15.943,0.032) -- (15.961,0.000) -- (15.997,0.000) -- (16.015,0.032) -- (16.034,0.000) -- (16.070,0.000);
\draw[cExact,thick,line join=round] (13.120,0.000) -- (13.156,0.000) -- (13.175,0.032) -- (13.193,0.000) -- (13.229,0.000) -- (13.247,0.032) -- (13.229,0.063) -- (13.266,0.063) -- (13.284,0.095) -- (13.302,0.063) -- (13.339,0.063) -- (13.320,0.032) -- (13.339,0.000) -- (13.375,0.000) -- (13.393,0.032) -- (13.411,0.000) -- (13.448,0.000) -- (13.466,0.032) -- (13.448,0.063) -- (13.484,0.063) -- (13.502,0.095) -- (13.484,0.126) -- (13.448,0.126) -- (13.466,0.158) -- (13.448,0.189) -- (13.484,0.189) -- (13.502,0.221) -- (13.521,0.189) -- (13.557,0.189) -- (13.575,0.221) -- (13.557,0.252) -- (13.593,0.252) -- (13.612,0.284) -- (13.630,0.252) -- (13.666,0.252) -- (13.648,0.221) -- (13.666,0.189) -- (13.703,0.189) -- (13.721,0.221) -- (13.739,0.189) -- (13.776,0.189) -- (13.757,0.158) -- (13.776,0.126) -- (13.739,0.126) -- (13.721,0.095) -- (13.739,0.063) -- (13.776,0.063) -- (13.757,0.032) -- (13.776,0.000) -- (13.812,0.000) -- (13.830,0.032) -- (13.848,0.000) -- (13.885,0.000) -- (13.903,0.032) -- (13.885,0.063) -- (13.921,0.063) -- (13.939,0.095) -- (13.958,0.063) -- (13.994,0.063) -- (13.976,0.032) -- (13.994,0.000) -- (14.030,0.000) -- (14.049,0.032) -- (14.067,0.000) -- (14.103,0.000);
\fill (13.120,0) circle (1.2pt);
\node[cExact,anchor=south] at (13.612,0.42) {$R=L/3$};
\node[anchor=north] at (14.595,-0.1) {$k=4$};
\end{tikzpicture}}

\newcommand{\ElevenFigCantor}{%
\begin{tikzpicture}[font=\scriptsize,baseline=(current bounding box.center)]
\draw[cWalk,line width=2.6pt] (0.000,0.000) -- (6.300,0.000);
\node[anchor=east] at (-0.1,0.000) {$k=0$};
\node[anchor=west] at (6.400,0.000) {1};
\draw[cWalk,line width=2.6pt] (0.000,-0.340) -- (2.100,-0.340) (4.200,-0.340) -- (6.300,-0.340);
\node[anchor=east] at (-0.1,-0.340) {$k=1$};
\node[anchor=west] at (6.400,-0.340) {2};
\draw[cWalk,line width=2.6pt] (0.000,-0.680) -- (0.700,-0.680) (1.400,-0.680) -- (2.100,-0.680) (4.200,-0.680) -- (4.900,-0.680) (5.600,-0.680) -- (6.300,-0.680);
\node[anchor=east] at (-0.1,-0.680) {$k=2$};
\node[anchor=west] at (6.400,-0.680) {4};
\draw[cWalk,line width=2.6pt] (0.000,-1.020) -- (0.233,-1.020) (0.467,-1.020) -- (0.700,-1.020) (1.400,-1.020) -- (1.633,-1.020) (1.867,-1.020) -- (2.100,-1.020) (4.200,-1.020) -- (4.433,-1.020) (4.667,-1.020) -- (4.900,-1.020) (5.600,-1.020) -- (5.833,-1.020) (6.067,-1.020) -- (6.300,-1.020);
\node[anchor=east] at (-0.1,-1.020) {$k=3$};
\node[anchor=west] at (6.400,-1.020) {8};
\draw[cWalk,line width=2.6pt] (0.000,-1.360) -- (0.078,-1.360) (0.156,-1.360) -- (0.233,-1.360) (0.467,-1.360) -- (0.544,-1.360) (0.622,-1.360) -- (0.700,-1.360) (1.400,-1.360) -- (1.478,-1.360) (1.556,-1.360) -- (1.633,-1.360) (1.867,-1.360) -- (1.944,-1.360) (2.022,-1.360) -- (2.100,-1.360) (4.200,-1.360) -- (4.278,-1.360) (4.356,-1.360) -- (4.433,-1.360) (4.667,-1.360) -- (4.744,-1.360) (4.822,-1.360) -- (4.900,-1.360) (5.600,-1.360) -- (5.678,-1.360) (5.756,-1.360) -- (5.833,-1.360) (6.067,-1.360) -- (6.144,-1.360) (6.222,-1.360) -- (6.300,-1.360);
\node[anchor=east] at (-0.1,-1.360) {$k=4$};
\node[anchor=west] at (6.400,-1.360) {16};
\draw[cWalk,line width=2.6pt] (0.000,-1.700) -- (0.026,-1.700) (0.052,-1.700) -- (0.078,-1.700) (0.156,-1.700) -- (0.181,-1.700) (0.207,-1.700) -- (0.233,-1.700) (0.467,-1.700) -- (0.493,-1.700) (0.519,-1.700) -- (0.544,-1.700) (0.622,-1.700) -- (0.648,-1.700) (0.674,-1.700) -- (0.700,-1.700) (1.400,-1.700) -- (1.426,-1.700) (1.452,-1.700) -- (1.478,-1.700) (1.556,-1.700) -- (1.581,-1.700) (1.607,-1.700) -- (1.633,-1.700) (1.867,-1.700) -- (1.893,-1.700) (1.919,-1.700) -- (1.944,-1.700) (2.022,-1.700) -- (2.048,-1.700) (2.074,-1.700) -- (2.100,-1.700) (4.200,-1.700) -- (4.226,-1.700) (4.252,-1.700) -- (4.278,-1.700) (4.356,-1.700) -- (4.381,-1.700) (4.407,-1.700) -- (4.433,-1.700) (4.667,-1.700) -- (4.693,-1.700) (4.719,-1.700) -- (4.744,-1.700) (4.822,-1.700) -- (4.848,-1.700) (4.874,-1.700) -- (4.900,-1.700) (5.600,-1.700) -- (5.626,-1.700) (5.652,-1.700) -- (5.678,-1.700) (5.756,-1.700) -- (5.781,-1.700) (5.807,-1.700) -- (5.833,-1.700) (6.067,-1.700) -- (6.093,-1.700) (6.119,-1.700) -- (6.144,-1.700) (6.222,-1.700) -- (6.248,-1.700) (6.274,-1.700) -- (6.300,-1.700);
\node[anchor=east] at (-0.1,-1.700) {$k=5$};
\node[anchor=west] at (6.400,-1.700) {32};
\end{tikzpicture}}

\newcommand{\ElevenTabFractal}{%
\begin{tabular}{@{}c*{6}{r}@{}}
\toprule
 & \multicolumn{2}{c}{Koch} & \multicolumn{2}{c}{Cantor} & \multicolumn{2}{c}{$n(R=3^{-k})$}\\
$k$ & 段數 & 總長 & 段數 & 總長 & Koch & Cantor\\
\midrule
0 & 1 & 1.0000 & 1 & 1.0000 & 1.0000 & 1.0000\\
1 & 4 & 1.3333 & 2 & 0.6667 & 0.2500 & 0.5000\\
2 & 16 & 1.7778 & 4 & 0.4444 & 0.0625 & 0.2500\\
3 & 64 & 2.3704 & 8 & 0.2963 & 0.0156 & 0.1250\\
4 & 256 & 3.1605 & 16 & 0.1975 & 0.0039 & 0.0625\\
5 & 1024 & 4.2140 & 32 & 0.1317 &  & \\
\midrule
斜率 & & $\to\infty$ & & $\to0$ & 1.2619 & 0.6309\\
\bottomrule
\end{tabular}}

\newcommand{\ElevenFigCollapse}{%
\begin{tikzpicture}[font=\scriptsize,>=Stealth,baseline=(current bounding box.center)]
\draw[->,gray] (0.000,0.000) -- (5.250,0.000) node[right,black] {$h$};
\draw[->,gray] (0.000,0.000) -- (0.000,2.375) node[above,black] {$m$};
\draw[gray] (0.000,0.000) -- ++(0,-0.06) node[below,black] {0};
\draw[gray] (1.250,0.000) -- ++(0,-0.06) node[below,black] {0.5};
\draw[gray] (2.500,0.000) -- ++(0,-0.06) node[below,black] {1};
\draw[gray] (3.750,0.000) -- ++(0,-0.06) node[below,black] {1.5};
\draw[gray] (5.000,0.000) -- ++(0,-0.06) node[below,black] {2};
\draw[gray] (0.000,0.850) -- ++(-0.06,0) node[left,black] {0.5};
\draw[gray] (0.000,1.700) -- ++(-0.06,0) node[left,black] {1};
\draw[cBar,thick] (0.000,0.000) -- (0.128,0.308) -- (0.256,0.512) -- (0.385,0.656) -- (0.513,0.768) -- (0.641,0.861) -- (0.769,0.941) -- (0.897,1.011) -- (1.026,1.074) -- (1.154,1.132) -- (1.282,1.185) -- (1.410,1.234) -- (1.538,1.280) -- (1.667,1.324) -- (1.795,1.365) -- (1.923,1.404) -- (2.051,1.441) -- (2.179,1.476) -- (2.308,1.510) -- (2.436,1.543) -- (2.564,1.574) -- (2.692,1.605) -- (2.821,1.634) -- (2.949,1.662) -- (3.077,1.690) -- (3.205,1.717) -- (3.333,1.743) -- (3.462,1.768) -- (3.590,1.793) -- (3.718,1.817) -- (3.846,1.840) -- (3.974,1.863) -- (4.103,1.885) -- (4.231,1.907) -- (4.359,1.929) -- (4.487,1.950) -- (4.615,1.970) -- (4.744,1.990) -- (4.872,2.010) -- (5.000,2.030);
\draw[cWalk,thick,densely dashed] (0.000,0.000) -- (0.128,0.171) -- (0.256,0.325) -- (0.385,0.457) -- (0.513,0.570) -- (0.641,0.667) -- (0.769,0.752) -- (0.897,0.828) -- (1.026,0.896) -- (1.154,0.959) -- (1.282,1.017) -- (1.410,1.070) -- (1.538,1.120) -- (1.667,1.167) -- (1.795,1.211) -- (1.923,1.253) -- (2.051,1.293) -- (2.179,1.331) -- (2.308,1.368) -- (2.436,1.403) -- (2.564,1.436) -- (2.692,1.469) -- (2.821,1.500) -- (2.949,1.530) -- (3.077,1.560) -- (3.205,1.588) -- (3.333,1.616) -- (3.462,1.642) -- (3.590,1.668) -- (3.718,1.694) -- (3.846,1.719) -- (3.974,1.743) -- (4.103,1.766) -- (4.231,1.789) -- (4.359,1.812) -- (4.487,1.834) -- (4.615,1.856) -- (4.744,1.877) -- (4.872,1.898) -- (5.000,1.918);
\draw[cExact,thick,densely dotted] (0.000,0.000) -- (0.128,0.087) -- (0.256,0.173) -- (0.385,0.256) -- (0.513,0.336) -- (0.641,0.412) -- (0.769,0.484) -- (0.897,0.552) -- (1.026,0.616) -- (1.154,0.677) -- (1.282,0.735) -- (1.410,0.789) -- (1.538,0.841) -- (1.667,0.890) -- (1.795,0.936) -- (1.923,0.981) -- (2.051,1.024) -- (2.179,1.065) -- (2.308,1.104) -- (2.436,1.142) -- (2.564,1.178) -- (2.692,1.213) -- (2.821,1.247) -- (2.949,1.280) -- (3.077,1.312) -- (3.205,1.342) -- (3.333,1.372) -- (3.462,1.401) -- (3.590,1.430) -- (3.718,1.457) -- (3.846,1.484) -- (3.974,1.510) -- (4.103,1.536) -- (4.231,1.561) -- (4.359,1.585) -- (4.487,1.609) -- (4.615,1.633) -- (4.744,1.656) -- (4.872,1.678) -- (5.000,1.700);
\fill[black] (0.625,0.850) circle (1.3pt);
\fill[black] (1.768,1.202) circle (1.3pt);
\fill[black] (5.000,1.700) circle (1.3pt);
\node[cBar,anchor=west] at (3.500,0.816) {$r=0.25$};
\node[cWalk,anchor=west] at (3.500,0.544) {$r=0.5$};
\node[cExact,anchor=west] at (3.500,0.272) {$r=1$};
\begin{scope}[xshift=6.3cm]
\draw[->,gray] (0.000,0.000) -- (5.250,0.000) node[right,black] {$x=h/r^{3/2}$};
\draw[->,gray] (0.000,0.000) -- (0.000,2.280) node[above,black] {$y=m/r^{1/2}$};
\draw[gray] (0.000,0.000) -- ++(0,-0.06) node[below,black] {0};
\draw[gray] (1.250,0.000) -- ++(0,-0.06) node[below,black] {1};
\draw[gray] (2.500,0.000) -- ++(0,-0.06) node[below,black] {2};
\draw[gray] (3.750,0.000) -- ++(0,-0.06) node[below,black] {3};
\draw[gray] (5.000,0.000) -- ++(0,-0.06) node[below,black] {4};
\draw[gray] (0.000,0.000) -- ++(-0.06,0) node[left,black] {0};
\draw[gray] (0.000,0.725) -- ++(-0.06,0) node[left,black] {0.5};
\draw[gray] (0.000,1.450) -- ++(-0.06,0) node[left,black] {1};
\draw[cBar,thick] (0.000,0.000) -- (0.128,0.147) -- (0.256,0.286) -- (0.385,0.413) -- (0.513,0.526) -- (0.641,0.627) -- (0.769,0.717) -- (0.897,0.799) -- (1.026,0.873) -- (1.154,0.942) -- (1.282,1.005) -- (1.410,1.064) -- (1.538,1.119) -- (1.667,1.171) -- (1.795,1.219) -- (1.923,1.266) -- (2.051,1.310) -- (2.179,1.352) -- (2.308,1.393) -- (2.436,1.431) -- (2.564,1.468) -- (2.692,1.504) -- (2.821,1.539) -- (2.949,1.572) -- (3.077,1.605) -- (3.205,1.636) -- (3.333,1.666) -- (3.462,1.696) -- (3.590,1.725) -- (3.718,1.753) -- (3.846,1.780) -- (3.974,1.806) -- (4.103,1.832) -- (4.231,1.858) -- (4.359,1.883) -- (4.487,1.907) -- (4.615,1.931) -- (4.744,1.954) -- (4.872,1.977) -- (5.000,1.999);
\draw[cWalk,thick,densely dashed] (0.000,0.000) -- (0.128,0.147) -- (0.256,0.286) -- (0.385,0.413) -- (0.513,0.526) -- (0.641,0.627) -- (0.769,0.717) -- (0.897,0.799) -- (1.026,0.873) -- (1.154,0.942) -- (1.282,1.005) -- (1.410,1.064) -- (1.538,1.119) -- (1.667,1.171) -- (1.795,1.219) -- (1.923,1.266) -- (2.051,1.310) -- (2.179,1.352) -- (2.308,1.393) -- (2.436,1.431) -- (2.564,1.468) -- (2.692,1.504) -- (2.821,1.539) -- (2.949,1.572) -- (3.077,1.605) -- (3.205,1.636) -- (3.333,1.666) -- (3.462,1.696) -- (3.590,1.725) -- (3.718,1.753) -- (3.846,1.780) -- (3.974,1.806) -- (4.103,1.832) -- (4.231,1.858) -- (4.359,1.883) -- (4.487,1.907) -- (4.615,1.931) -- (4.744,1.954) -- (4.872,1.977) -- (5.000,1.999);
\draw[cExact,thick,densely dotted] (0.000,0.000) -- (0.128,0.147) -- (0.256,0.286) -- (0.385,0.413) -- (0.513,0.526) -- (0.641,0.627) -- (0.769,0.717) -- (0.897,0.799) -- (1.026,0.873) -- (1.154,0.942) -- (1.282,1.005) -- (1.410,1.064) -- (1.538,1.119) -- (1.667,1.171) -- (1.795,1.219) -- (1.923,1.266) -- (2.051,1.310) -- (2.179,1.352) -- (2.308,1.393) -- (2.436,1.431) -- (2.564,1.468) -- (2.692,1.504) -- (2.821,1.539) -- (2.949,1.572) -- (3.077,1.605) -- (3.205,1.636) -- (3.333,1.666) -- (3.462,1.696) -- (3.590,1.725) -- (3.718,1.753) -- (3.846,1.780) -- (3.974,1.806) -- (4.103,1.832) -- (4.231,1.858) -- (4.359,1.883) -- (4.487,1.907) -- (4.615,1.931) -- (4.744,1.954) -- (4.872,1.977) -- (5.000,1.999);
\draw[gray,densely dotted] (2.500,0) -- (2.500,1.450) -- (0,1.450);
\fill[black] (2.500,1.450) circle (1.5pt);
\node[anchor=west,align=left] at (2.625,0.899) {三個黑點\\都到 $(2,1)$};
\node[anchor=west] at (3.750,2.175) {$y+y^3=x$};
\end{scope}
\end{tikzpicture}}

\newcommand{\ElevenFigChain}{%
\begin{tikzpicture}[font=\scriptsize,>=Stealth,baseline=(current bounding box.center)]
\draw[->,gray] (0.000,0.000) -- (5.350,0.000) node[right,black] {$r/a$};
\draw[->,gray] (0.000,0.000) -- (0.000,3.150) node[above,black] {$r''/a$};
\draw[gray] (0.000,0.000) -- ++(0,-0.06) node[below,black] {0};
\draw[gray] (1.275,0.000) -- ++(0,-0.06) node[below,black] {0.5};
\draw[gray] (2.550,0.000) -- ++(0,-0.06) node[below,black] {1};
\draw[gray] (3.825,0.000) -- ++(0,-0.06) node[below,black] {1.5};
\draw[gray] (5.100,0.000) -- ++(0,-0.06) node[below,black] {2};
\draw[gray] (0.000,0.580) -- ++(-0.06,0) node[left,black] {2};
\draw[gray] (0.000,1.160) -- ++(-0.06,0) node[left,black] {4};
\draw[gray] (0.000,1.740) -- ++(-0.06,0) node[left,black] {6};
\draw[gray] (0.000,2.320) -- ++(-0.06,0) node[left,black] {8};
\draw[gray] (0.000,2.900) -- ++(-0.06,0) node[left,black] {10};
\draw[gray,densely dashed] (0,0) -- (5.100,0.580) node[above,black] {$r''=r$};
\draw[cBar,thick] (0.000,0.000) -- (0.114,0.052) -- (0.228,0.106) -- (0.342,0.161) -- (0.456,0.216) -- (0.569,0.273) -- (0.683,0.332) -- (0.797,0.391) -- (0.911,0.451) -- (1.025,0.513) -- (1.139,0.576) -- (1.253,0.640) -- (1.367,0.705) -- (1.480,0.771) -- (1.594,0.839) -- (1.708,0.907) -- (1.822,0.977) -- (1.936,1.048) -- (2.050,1.120) -- (2.164,1.193) -- (2.278,1.267) -- (2.391,1.343) -- (2.505,1.420) -- (2.619,1.497) -- (2.733,1.576) -- (2.847,1.657) -- (2.961,1.738) -- (3.075,1.820) -- (3.189,1.904) -- (3.302,1.989) -- (3.416,2.075) -- (3.530,2.162) -- (3.644,2.250) -- (3.758,2.339) -- (3.872,2.430) -- (3.986,2.522) -- (4.100,2.614) -- (4.213,2.708) -- (4.327,2.804) -- (4.441,2.900);
\node[cBar,anchor=east] at (4.391,2.900) {$r''=r(r+4a)/a$};
\draw[cWalk,thick,->] (0.128,0.000) -- (0.128,0.059) -- (0.516,0.059) -- (0.516,0.247) -- (2.170,0.247) -- (2.170,1.197);
\fill[cWalk] (0.128,0.059) circle (1.1pt);
\fill[cWalk] (0.516,0.247) circle (1.1pt);
\fill[cWalk] (2.170,1.197) circle (1.1pt);
\node[cWalk,anchor=west] at (2.220,1.197) {$r_3=4.1282$};
\node[cWalk,anchor=west,align=left] at (0.153,2.088) {$r_0=0.05\to0.2025$\\$\to0.8510\to4.1282$};
\fill (0,0) circle (1.4pt);
\end{tikzpicture}}

\newcommand{\ElevenTabSmallR}{%
\begin{tabular}{@{}crrr@{}}
\toprule
$k$ & $\dfrac{r^{(k)}}{4^kr}$ & $\dfrac{h^{(k)}}{2^{3k/2}h}$ & $\dfrac{\text{自旋因子}}{2^{k/2}}$\\
\midrule
1 & 1.0000 & 1.0000 & 1.0000\\
2 & 1.0001 & 1.0000 & 1.0001\\
3 & 1.0005 & 1.0000 & 1.0005\\
4 & 1.0021 & 1.0000 & 1.0021\\
5 & 1.0086 & 1.0000 & 1.0085\\
\bottomrule
\end{tabular}}

\newcommand{\ElevenFigXmap}{%
\begin{tikzpicture}[font=\scriptsize,>=Stealth,baseline=(current bounding box.center)]
\draw[->,gray] (0.000,0.000) -- (3.750,0.000) node[right,black] {$x=e^{-2K}$};
\draw[->,gray] (0.000,0.000) -- (0.000,3.750) node[above,black] {$x'$};
\draw[gray] (0.000,0.000) -- ++(0,-0.06) node[below,black] {0};
\draw[gray] (1.750,0.000) -- ++(0,-0.06) node[below,black] {0.5};
\draw[gray] (3.500,0.000) -- ++(0,-0.06) node[below,black] {1};
\draw[gray] (0.000,1.750) -- ++(-0.06,0) node[left,black] {0.5};
\draw[gray] (0.000,3.500) -- ++(-0.06,0) node[left,black] {1};
\draw[gray,densely dashed] (0,0) -- (3.500,3.500);
\draw[cBar,thick] (0.000,0.000) -- (0.090,0.179) -- (0.179,0.358) -- (0.269,0.535) -- (0.359,0.710) -- (0.449,0.883) -- (0.538,1.052) -- (0.628,1.217) -- (0.718,1.378) -- (0.808,1.534) -- (0.897,1.684) -- (0.987,1.829) -- (1.077,1.968) -- (1.167,2.100) -- (1.256,2.226) -- (1.346,2.345) -- (1.436,2.458) -- (1.526,2.564) -- (1.615,2.663) -- (1.705,2.756) -- (1.795,2.842) -- (1.885,2.922) -- (1.974,2.996) -- (2.064,3.063) -- (2.154,3.124) -- (2.244,3.180) -- (2.333,3.231) -- (2.423,3.276) -- (2.513,3.316) -- (2.603,3.352) -- (2.692,3.383) -- (2.782,3.410) -- (2.872,3.433) -- (2.962,3.452) -- (3.051,3.467) -- (3.141,3.480) -- (3.231,3.489) -- (3.321,3.495) -- (3.410,3.499) -- (3.500,3.500);
\draw[cWalk,thick,->] (0.474,0.000) -- (0.474,0.930) -- (0.930,0.930) -- (0.930,1.738) -- (1.738,1.738) -- (1.738,2.788) -- (2.788,2.788) -- (2.788,3.411) -- (3.411,3.411) -- (3.411,3.499);
\fill (0,0) circle (1.5pt);
\fill (3.500,3.500) circle (1.5pt);
\node[anchor=north west,align=left] at (0.08,-0.3) {$x=0$：$T=0$，斜率 2（不穩定）};
\node[anchor=west,align=left] at (3.580,3.500) {$x=1$：$T=\infty$\\斜率 0（穩定）};
\node[cBar,anchor=west] at (1.925,0.770) {藍：$x'=\frac{2x}{1+x^2}$};
\end{tikzpicture}}

\newcommand{\ElevenTabIsing}{%
\begin{tabular}{@{}rrr@{}}
\toprule
$N$ & 暴力＝轉移矩陣＝連乘 & 書上指數\\
\midrule
$4$ & 121.2329\,$\checkmark$ & {\color{cWrong}18.5498}\\
$8$ & 9160.4387\,$\checkmark$ & {\color{cWrong}143.2705}\\
$16$ & 68584536.5912\,$\checkmark$ & {\color{cWrong}11787.1848}\\
\bottomrule
\end{tabular}}

\newcommand{\ElevenFigShell}{%
\begin{tikzpicture}[font=\scriptsize,baseline=(current bounding box.center)]
\fill[cWalk!25,even odd rule] (-1.25,-1.25) rectangle (1.25,1.25) (-0.625,-0.625) rectangle (0.625,0.625);
\draw[cWalk,thick] (-1.25,-1.25) rectangle (1.25,1.25);
\draw[cBar,thick] (-0.625,-0.625) rectangle (0.625,0.625);
\draw[gray,->] (-1.45,0) -- (1.55,0) node[right,black] {$q_x$};
\draw[gray,->] (0,-1.45) -- (0,1.55) node[above,black] {$q_y$};
\node[anchor=north west] at (1.25,0) {$\frac\pi c$};
\node[anchor=north west,cBar] at (0.625,0) {$\frac\pi{2c}$};
\node[align=center] at (0,-1.8) {殼層（橙）被積掉：\\$1-2^{-d}=\frac{3}{4}$};
\end{tikzpicture}}

\newcommand{\ElevenFigFlow}{%
\begin{tikzpicture}[font=\scriptsize,>=Stealth,baseline=(current bounding box.center)]
\draw[->,gray] (-2.100,0.000) -- (2.350,0.000) node[right,black] {$r$};
\draw[->,gray] (0.000,0.000) -- (0.000,3.400) node[above,black] {$h$};
\draw[gray] (-2.100,0.000) -- ++(0,-0.06) node[below,black] {-1};
\draw[gray] (-1.050,0.000) -- ++(0,-0.06) node[below,black] {-0.5};
\draw[gray] (0.000,0.000) -- ++(0,-0.06) node[below,black] {0};
\draw[gray] (1.050,0.000) -- ++(0,-0.06) node[below,black] {0.5};
\draw[gray] (2.100,0.000) -- ++(0,-0.06) node[below,black] {1};
\draw[gray] (0.000,1.050) -- ++(-0.06,0) node[left,black] {0.5};
\draw[gray] (0.000,2.100) -- ++(-0.06,0) node[left,black] {1};
\draw[gray] (0.000,3.150) -- ++(-0.06,0) node[left,black] {1.5};
\draw[cBar,thick,postaction={decorate},decoration={markings,mark=at position 0.55 with {\arrow{Stealth[length=3pt]}}}] (0.021,0.021) -- (0.023,0.024) -- (0.026,0.027) -- (0.029,0.031) -- (0.032,0.035) -- (0.035,0.040) -- (0.039,0.045) -- (0.043,0.052) -- (0.048,0.059) -- (0.053,0.067) -- (0.059,0.076) -- (0.065,0.086) -- (0.072,0.098) -- (0.080,0.112) -- (0.089,0.127) -- (0.098,0.144) -- (0.109,0.164) -- (0.121,0.187) -- (0.134,0.212) -- (0.148,0.241) -- (0.164,0.274) -- (0.182,0.312) -- (0.201,0.355) -- (0.223,0.403) -- (0.247,0.459) -- (0.274,0.521) -- (0.304,0.593) -- (0.337,0.674) -- (0.373,0.767) -- (0.414,0.872) -- (0.459,0.991) -- (0.508,1.127) -- (0.563,1.282) -- (0.624,1.457) -- (0.692,1.657) -- (0.767,1.884) -- (0.850,2.142) -- (0.941,2.436) -- (1.043,2.770) -- (1.156,3.150);
\fill[cBar] (0.021,0.021) circle (0.9pt);
\draw[cBar,thick,postaction={decorate},decoration={markings,mark=at position 0.55 with {\arrow{Stealth[length=3pt]}}}] (-0.021,0.021) -- (-0.023,0.024) -- (-0.026,0.027) -- (-0.029,0.031) -- (-0.032,0.035) -- (-0.035,0.040) -- (-0.039,0.045) -- (-0.043,0.052) -- (-0.048,0.059) -- (-0.053,0.067) -- (-0.059,0.076) -- (-0.065,0.086) -- (-0.072,0.098) -- (-0.080,0.112) -- (-0.089,0.127) -- (-0.098,0.144) -- (-0.109,0.164) -- (-0.121,0.187) -- (-0.134,0.212) -- (-0.148,0.241) -- (-0.164,0.274) -- (-0.182,0.312) -- (-0.201,0.355) -- (-0.223,0.403) -- (-0.247,0.459) -- (-0.274,0.521) -- (-0.304,0.593) -- (-0.337,0.674) -- (-0.373,0.767) -- (-0.414,0.872) -- (-0.459,0.991) -- (-0.508,1.127) -- (-0.563,1.282) -- (-0.624,1.457) -- (-0.692,1.657) -- (-0.767,1.884) -- (-0.850,2.142) -- (-0.941,2.436) -- (-1.043,2.770) -- (-1.156,3.150);
\fill[cBar] (-0.021,0.021) circle (0.9pt);
\draw[cBar,thick,postaction={decorate},decoration={markings,mark=at position 0.55 with {\arrow{Stealth[length=3pt]}}}] (0.063,0.006) -- (0.069,0.007) -- (0.075,0.008) -- (0.083,0.009) -- (0.090,0.010) -- (0.099,0.011) -- (0.108,0.012) -- (0.118,0.014) -- (0.129,0.015) -- (0.142,0.017) -- (0.155,0.019) -- (0.169,0.022) -- (0.185,0.024) -- (0.203,0.027) -- (0.222,0.030) -- (0.243,0.034) -- (0.266,0.038) -- (0.291,0.043) -- (0.318,0.048) -- (0.348,0.053) -- (0.380,0.060) -- (0.416,0.067) -- (0.455,0.075) -- (0.498,0.084) -- (0.545,0.093) -- (0.596,0.105) -- (0.653,0.117) -- (0.714,0.131) -- (0.781,0.147) -- (0.855,0.164) -- (0.935,0.184) -- (1.023,0.205) -- (1.119,0.230) -- (1.224,0.257) -- (1.340,0.288) -- (1.466,0.322) -- (1.604,0.360) -- (1.754,0.403) -- (1.919,0.451) -- (2.100,0.505);
\fill[cBar] (0.063,0.006) circle (0.9pt);
\draw[cBar,thick,postaction={decorate},decoration={markings,mark=at position 0.55 with {\arrow{Stealth[length=3pt]}}}] (-0.063,0.006) -- (-0.069,0.007) -- (-0.075,0.008) -- (-0.083,0.009) -- (-0.090,0.010) -- (-0.099,0.011) -- (-0.108,0.012) -- (-0.118,0.014) -- (-0.129,0.015) -- (-0.142,0.017) -- (-0.155,0.019) -- (-0.169,0.022) -- (-0.185,0.024) -- (-0.203,0.027) -- (-0.222,0.030) -- (-0.243,0.034) -- (-0.266,0.038) -- (-0.291,0.043) -- (-0.318,0.048) -- (-0.348,0.053) -- (-0.380,0.060) -- (-0.416,0.067) -- (-0.455,0.075) -- (-0.498,0.084) -- (-0.545,0.093) -- (-0.596,0.105) -- (-0.653,0.117) -- (-0.714,0.131) -- (-0.781,0.147) -- (-0.855,0.164) -- (-0.935,0.184) -- (-1.023,0.205) -- (-1.119,0.230) -- (-1.224,0.257) -- (-1.340,0.288) -- (-1.466,0.322) -- (-1.604,0.360) -- (-1.754,0.403) -- (-1.919,0.451) -- (-2.100,0.505);
\fill[cBar] (-0.063,0.006) circle (0.9pt);
\fill[cWalk] (0.155,0.256) circle (1.4pt) node[right] {$\ell=1$};
\fill[cWalk] (1.147,3.117) circle (1.4pt) node[right] {$\ell=2$};
\node[anchor=west,align=left] at (1.302,2.100) {$h\propto|r|^{5/4}$\\都離開原點};
\fill (0,0) circle (1.5pt);
\end{tikzpicture}}

\newcommand{\ElevenTabAlpha}{%
\begin{tabular}{@{}rrrrr@{}}
\toprule
$k$ & $q=\frac{2\pi k}8$ & $a(1-\cos q)$ & $\frac12aq^2$ & 書上\\
\midrule
0 & 0 & 0 & 0 & {\color{cWrong}0}\\
1 & 0.7854 & 0.2929 & 0.3084 & {\color{cWrong}-0.5858}\\
2 & 1.5708 & 1.0000 & 1.2337 & {\color{cWrong}-2.0000}\\
3 & 2.3562 & 1.7071 & 2.7758 & {\color{cWrong}-3.4142}\\
4 & 3.1416 & 2.0000 & 4.9348 & {\color{cWrong}-4.0000}\\
\bottomrule
\end{tabular}}

\newcommand{\ElevenFigToy}{%
\begin{tikzpicture}[font=\scriptsize,>=Stealth,scale=0.85,baseline=(current bounding box.center)]
\draw[cExact,densely dashed] (-3.1,-1.55) -- (3.1,1.55) node[above left,align=right] {相關 $(2,1)$\\$\Lambda=2$};
\draw[cWrong,densely dashed] (-0.85,1.7) -- (1.1,-2.2) node[right,align=left] {無關 $(1,-2)$\\$\Lambda=-3$（臨界線）};
\draw[cBar,thick,postaction={decorate},decoration={markings,mark=at position 0.55 with {\arrow{Stealth[length=3pt]}}}] (1.000,-2.000) -- (0.902,-1.803) -- (0.813,-1.626) -- (0.733,-1.466) -- (0.661,-1.322) -- (0.596,-1.192) -- (0.538,-1.075) -- (0.485,-0.969) -- (0.437,-0.874) -- (0.394,-0.788) -- (0.355,-0.711) -- (0.320,-0.641) -- (0.289,-0.578) -- (0.261,-0.521) -- (0.235,-0.470) -- (0.212,-0.424) -- (0.191,-0.382) -- (0.172,-0.345) -- (0.155,-0.311) -- (0.140,-0.280) -- (0.126,-0.253) -- (0.114,-0.228) -- (0.103,-0.205) -- (0.093,-0.185) -- (0.084,-0.167) -- (0.075,-0.151) -- (0.068,-0.136) -- (0.061,-0.122) -- (0.055,-0.110) -- (0.050,-0.100);
\fill[cBar] (1.000,-2.000) circle (1pt);
\draw[cBar,thick,postaction={decorate},decoration={markings,mark=at position 0.55 with {\arrow{Stealth[length=3pt]}}}] (1.100,-1.950) -- (1.009,-1.750) -- (0.928,-1.569) -- (0.856,-1.405) -- (0.793,-1.256) -- (0.737,-1.122) -- (0.689,-1.000) -- (0.647,-0.888) -- (0.611,-0.787) -- (0.580,-0.695) -- (0.555,-0.611) -- (0.534,-0.534) -- (0.518,-0.464) -- (0.506,-0.399) -- (0.498,-0.339) -- (0.493,-0.283) -- (0.493,-0.231) -- (0.495,-0.183) -- (0.501,-0.138) -- (0.511,-0.095) -- (0.524,-0.054) -- (0.539,-0.015) -- (0.559,0.023) -- (0.581,0.059) -- (0.607,0.095) -- (0.636,0.130) -- (0.669,0.165) -- (0.705,0.199) -- (0.745,0.234) -- (0.789,0.270);
\fill[cBar] (1.100,-1.950) circle (1pt);
\draw[cBar,thick,postaction={decorate},decoration={markings,mark=at position 0.55 with {\arrow{Stealth[length=3pt]}}}] (0.900,-2.050) -- (0.795,-1.857) -- (0.698,-1.684) -- (0.610,-1.528) -- (0.529,-1.388) -- (0.455,-1.263) -- (0.386,-1.151) -- (0.323,-1.051) -- (0.263,-0.961) -- (0.208,-0.881) -- (0.156,-0.810) -- (0.107,-0.748) -- (0.060,-0.692) -- (0.015,-0.644) -- (-0.028,-0.601) -- (-0.069,-0.564) -- (-0.110,-0.533) -- (-0.151,-0.506) -- (-0.191,-0.484) -- (-0.231,-0.466) -- (-0.271,-0.451) -- (-0.312,-0.441) -- (-0.353,-0.433) -- (-0.396,-0.429) -- (-0.440,-0.429) -- (-0.485,-0.431) -- (-0.533,-0.436) -- (-0.582,-0.444) -- (-0.634,-0.455) -- (-0.689,-0.469);
\fill[cBar] (0.900,-2.050) circle (1pt);
\draw[cBar,thick,postaction={decorate},decoration={markings,mark=at position 0.55 with {\arrow{Stealth[length=3pt]}}}] (0.900,-0.800) -- (0.879,-0.687) -- (0.866,-0.584) -- (0.859,-0.487) -- (0.858,-0.398) -- (0.863,-0.314) -- (0.874,-0.235) -- (0.891,-0.161) -- (0.913,-0.090) -- (0.941,-0.022) -- (0.975,0.043) -- (1.014,0.107) -- (1.060,0.169) -- (1.111,0.230) -- (1.168,0.290) -- (1.231,0.351) -- (1.301,0.412) -- (1.378,0.474) -- (1.462,0.537) -- (1.553,0.601) -- (1.652,0.668) -- (1.759,0.737) -- (1.875,0.809) -- (2.000,0.884) -- (2.135,0.963) -- (2.281,1.046) -- (2.437,1.134) -- (2.605,1.226) -- (2.786,1.324) -- (2.981,1.428);
\fill[cBar] (0.900,-0.800) circle (1pt);
\draw[cBar,thick,postaction={decorate},decoration={markings,mark=at position 0.55 with {\arrow{Stealth[length=3pt]}}}] (0.100,-1.200) -- (0.022,-1.116) -- (-0.053,-1.043) -- (-0.125,-0.979) -- (-0.196,-0.925) -- (-0.267,-0.879) -- (-0.336,-0.840) -- (-0.406,-0.809) -- (-0.476,-0.784) -- (-0.547,-0.766) -- (-0.620,-0.754) -- (-0.694,-0.748) -- (-0.771,-0.747) -- (-0.850,-0.751) -- (-0.933,-0.760) -- (-1.020,-0.775) -- (-1.110,-0.794) -- (-1.206,-0.818) -- (-1.306,-0.847) -- (-1.413,-0.882) -- (-1.526,-0.921) -- (-1.645,-0.965) -- (-1.772,-1.015) -- (-1.908,-1.070) -- (-2.052,-1.130) -- (-2.205,-1.197) -- (-2.369,-1.270) -- (-2.544,-1.349) -- (-2.731,-1.435) -- (-2.931,-1.528);
\fill[cBar] (0.100,-1.200) circle (1pt);
\fill (0,0) circle (1.8pt) node[left=2pt] {不動點};
\end{tikzpicture}}

\newcommand{\ElevenTabExp}{%
\begin{tabular}{@{}lcccc@{}}
\toprule
 & \shortstack{Gaussian\\$d=3$} & \shortstack{Gaussian\\$d=4$} & \shortstack{二維\\Ising} & \shortstack{二維\\自迴避}\\
\midrule
$\Lambda_r$ & $2$ & $2$ & $1$ & $\frac{4}{3}$\\
$\Lambda_h$ & $\frac{5}{2}$ & $3$ & $\frac{15}{8}$ & $\frac{91}{48}$\\
\midrule
$\alpha=2-\frac d{\Lambda_r}$ & $\frac{1}{2}$ & $0$ & $0$ & $\frac{1}{2}$\\
$\beta=\frac{d-\Lambda_h}{\Lambda_r}$ & $\frac{1}{4}$ & $\frac{1}{2}$ & $\frac{1}{8}$ & $\frac{5}{64}$\\
$\gamma=\frac{2\Lambda_h-d}{\Lambda_r}$ & $1$ & $1$ & $\frac{7}{4}$ & $\frac{43}{32}$\\
$\delta=\frac{\Lambda_h}{d-\Lambda_h}$ & $5$ & $3$ & $15$ & $\frac{91}{5}$\\
$\nu=\frac1{\Lambda_r}$ & $\frac{1}{2}$ & $\frac{1}{2}$ & $1$ & $\frac{3}{4}$\\
\midrule
$\alpha+2\beta+\gamma$ & 2 & 2 & 2 & 2\\
＝第 9 章 & & $\checkmark$ & $\checkmark$ & $\checkmark$\\
\bottomrule
\end{tabular}}

\newcommand{\ElevenTabErrata}{%
\begin{tabular}{@{}lll@{}}
\toprule
位置 & 書上 & 應為\\
\midrule
p.~263，(11.16) 後 & {\color{cWrong}$f(h/r^{3/2})\to h^{1/3}r$} & {\color{cExact}$f(h/r^{3/2})\to h^{1/3}r^{-1/2}$}\\
p.~265，(11.31) & {\color{cWrong}$\big(\frac{2\pi}{r+2a}\big)^{N/2}$} & {\color{cExact}$\big(\frac{2\pi}{r+2a}\big)^{N/4}$}\\
pp.~267--268，Fig.~11.6--11.7 & {\color{cWrong}軸 $r/\sqrt a$、$r''/\sqrt a$} & {\color{cExact}$r/a$、$r''/a$}\\
p.~267，正文 & {\color{cWrong}the reduced temperature, $t''$} & {\color{cExact}$r''$}\\
p.~269，(11.46)(11.47) & {\color{cWrong}$m=2^{-n/2}f(r2^{2n},h2^{3n/2})$} & {\color{cExact}$m=2^{+n/2}f(r2^{2n},h2^{3n/2})$}\\
p.~270，(11.53) & {\color{cWrong}$\cosh K(S_2+S_4)=2\cosh2K$（$S_2=S_4$）} & {\color{cExact}$2\cosh K(S_2+S_4)=2\cosh2K$}\\
p.~271，(11.63) & {\color{cWrong}$\cdots A(K_j)^{N/2^{j+1}}\cdots$} & {\color{cExact}$\cdots A(K_j)^{N/2^j}\cdots$}\\
p.~273，(11.66)(11.67) & {\color{cWrong}$\alpha(q)=-2a(1-\cos qc)\sim ac^2q^2$} & {\color{cExact}$\alpha(q)=a(1-\cos qc)\sim\frac12ac^2q^2$}\\
\bottomrule
\end{tabular}}


\section{標度、碎形與重整化（第 11 章）}

\begin{keybox}[先看這裡]
第 8、9 章說：臨界點附近的系統\textbf{沒有自己的尺度}（關聯長度 $\xi=\infty$），放大之後看起來還一樣——就像「跳蚤身上還有更小的跳蚤，一層又一層」。這一章把這句話變成計算工具：
先用最簡單的\textbf{碎形}（Koch 曲線、Cantor 集）說明「放大 $\mathcal M$ 倍、份量變 $\mathcal R$ 倍」怎麼給出非整數的維度 $\ln\mathcal R/\ln\mathcal M$；
再把系統放大、把一半的自由度\textbf{積掉}（抽格，decimation），看模型的參數怎麼變——參數的\textbf{遞迴關係}與\textbf{流}；
流停下來的地方是\textbf{不動點}，在不動點附近把流線性化，特徵值 $\Lambda$ 決定一切：$\Lambda>0$ 的方向\textbf{相關}（會被放大，要調到臨界才不見），$\Lambda<0$ 的方向\textbf{無關}（細節被忘掉——普適性），而 $\nu=1/\Lambda_r$。

\textbf{記號}：碎形維度 $p$（$=d_f$）；Landau 的約化溫度 $r$、外場 $h$、磁化 $m$（本系列照第 9 章寫 $\frac r2m^2+\frac u4m^4-hm$，$u=1$）；一維 Gaussian 鏈的耦合 $a$（$=1$）；
Ising 的 $K=J/k_BT$、外場 $H=h/k_BT$；抽格 $b=2$、連續重整化 $b=e^\ell$；相關特徵值 $\Lambda_r,\Lambda_h$。
\begin{center}
\small
\begin{tabular}{@{}lll@{\qquad}lll@{}}
\toprule
節 & 內容 & 動畫 & 節 & 內容 & 動畫\\
\midrule
11.1.1 & 碎形、維度 & W1101 & 11.2 例題 & 一維 Ising 的抽格 & W1104\\
11.1.2 & 被動的重整化 & W1102 & 11.2.1--11.2.2 & 動量殼、連續重整化 & W1105\\
11.2 前半 & 一維 Gaussian 鏈 & W1103 & 11.3 & 不動點、臨界指數 & W1106\\
\bottomrule
\end{tabular}
\end{center}
數字一律由 \texttt{codes/utils/walk\_rg.py} 重算，逐式校對見 \texttt{ref/md/ch11\_formulas.md}，
圖表由 \texttt{ch11/make\_figs.py} 產生並先做數字檢查。連續播放：\texttt{Ch11Demo}（W1101--W1106）。
\end{keybox}

% =====================================================================
\subsection{碎形：放大之後還一樣}
\leavevmode\bk{pp.~255--261}\ \anim{W1101}\par\nopagebreak

\textbf{問}：什麼東西放大之後還長得一樣？

\textbf{直觀}：點和直線放大後不變（平凡）；圓弧越放大越直——圓有自己的尺度（半徑），放大到比半徑小的範圍就看不出是圓。
\textbf{Koch 曲線}（圖\ref{eleven-fig-fractal} 上）沒有這種尺度：每一代把每段的中間三分之一換成等邊三角形的兩邊，一段變 4 段、每段長 $\frac13$。
第 $k$ 代 $4^k$ 段、總長 $(\frac43)^k\to\infty$（第 5 代已是 4.2140）。取第 5 代的前四分之一放大 3 倍，和整條疊在一起一模一樣（自相似）。
\textbf{Cantor 集}（圖\ref{eleven-fig-fractal} 左下）反過來：每段挖掉中間三分之一，$2^k$ 段、總長 $(\frac23)^k\to0$。一個無限長、一個長度是 0——「長度」這把尺對它們不適用，要換一把。

\textbf{質量－半徑}（Fig.~11.5，(11.1)--(11.5)）：以物體上一點為圓心、半徑 $R$ 的圓內有多少「份量」$n(R)$？直線上 $n=2R\propto R^1$，平面上 $n=\pi R^2\propto R^2$：一般寫 $n(R)=KR^p$，指數 $p$ 就是維度。
Koch 曲線（總份量算 1）以左端為圓心：半徑 $L$ 的圓包住整條，$n(L)=1$；半徑 $L/3$ 的圓裡正好是前四分之一（圖 \ref{eleven-fig-fractal} 上右的綠色），而前四分之一放大 3 倍就是整條，所以每縮 3 倍都一樣：
\[n(R/3)=\tfrac14\,n(R)\ \Rightarrow\ K\big(\tfrac R3\big)^p=\tfrac14KR^p\ \Rightarrow\ 3^{-p}=\tfrac14\ \Rightarrow\ p=\frac{\ln4}{\ln3}=1.2619 .\tag{11.1--11.5}\]
\textbf{算到底}（圖 \ref{eleven-fig-fractal} 右下的表）：$R=3^{-k}$ 時 $n=4^{-k}$（0.2500、0.0625……），斜率正好 1.2619；Cantor 集 $n=2^{-k}$，斜率 $\frac{\ln2}{\ln3}=0.6309$（練習 11.2）。
〔注〕$n(L/3)=\frac14n(L)$ 是精確的：整條 Koch 曲線落在以左端為頂點、張角 $30^\circ$ 的錐裡，第 2 段（轉 $60^\circ$）以後的點離左端都 $\ge L/3$。
改取 $R=2^{-j}$（$j=0..8$）斜率是 1.2571：$n(R)$ 是 $R^p$ 乘上以 $\ln3$ 為週期的小抖動。
\textbf{第二種量法：盒子計數}。用邊長 $\varepsilon$ 的方格蓋住曲線，數有曲線的格子 $N(\varepsilon)\propto\varepsilon^{-p}$：$\varepsilon=\frac19,\frac1{27},\frac1{81},\frac1{243}$ 得 22、83、325、1291 個（每縮 3 倍約乘 4），
斜率 1.2362；Cantor 集 0.6254。有限的代數讓它略偏低，但都指向 $\frac{\ln4}{\ln3}$、$\frac{\ln2}{\ln3}$。

\textbf{複本公式}（(11.6)）：若物體由 $\mathcal R$ 個小複本組成、每個放大 $\mathcal M$ 倍就變回原物，則 $\mathcal M^p=\mathcal R$：
\[p=\frac{\ln\mathcal R}{\ln\mathcal M}.\tag{11.6}\]
線段（3 份、放大 3 倍）1.0000；正方形（9 份、3 倍）2.0000（練習 11.1）；Koch $(4,3)$ 1.2619；Cantor $(2,3)$ 0.6309（練習 11.3）；Sierpinski 三角形 $(3,2)$ 1.5850（本系列加的）。Koch 曲線介於線和面之間。

\textbf{漫步也是碎形}（p.~261）：普通漫步 $N$ 步走 $R\sim N^{1/2}$，把步數當份量就是 $N\propto R^2$：\textbf{任何維度裡普通漫步的碎形維度都是 2}。
\textbf{核對}（本系列）：迴轉半徑 $\langle R_g^2\rangle=\frac{N(N+2)}{6(N+1)}$（每步長 1、任何維度；$N\le5$ 與 $4^N$ 條全部列舉相同）。400 條二維漫步（種子 1101）在 $N=64,256,1024,4096$ 得
11.23、44.65、163.53、696.53（公式 10.83、42.83、170.83、682.83），$\ln N$ 對 $\ln R_g$ 的斜率 2.0258（公式 2.0069）$\to2$。
\textbf{自迴避漫步} $R\sim N^\nu$，所以 $d_f=1/\nu$：二維 $\frac43$（第 7 章精確列舉的有效值 $N=10,12,14$：1.3848、1.3788、1.3740 $\to\frac43$）、三維 $1/0.588=1.70$（引用）、四維以上 2——不能走回頭路，漫步比較「稀」。

\begin{figure}[htbp]
\centering
\ElevenFigKoch\\[6pt]
\begin{minipage}[c]{0.5\linewidth}\centering\ElevenFigCantor\end{minipage}\hfill
\begin{minipage}[c]{0.48\linewidth}\centering\footnotesize\setlength{\tabcolsep}{3pt}\ElevenTabFractal\end{minipage}
\caption{碎形。上：Koch 曲線第 0--4 代；半徑 $L/3$（綠虛線）的圓裡正好是前四分之一（綠）。左下：Cantor 集。右下：段數、總長與 $n(R)$（最後一列是斜率）}\label{eleven-fig-fractal}
\end{figure}

\textbf{揭曉}：碎形維度＝「放大 $\mathcal M$ 倍、份量變 $\mathcal R$ 倍」的對數比；重整化群就是把「放大之後還一樣」這個想法用在漲落的系統上。

% =====================================================================
\subsection{被動的重整化：平均場與 Gaussian 模型}
\leavevmode\bk{pp.~261--264}\ \anim{W1102}\par\nopagebreak

\textbf{問}：把整個磁鐵放大 $b$ 倍，自由能的每一項會怎樣？

\textbf{Landau}（(11.7)--(11.10)）：平均場裡磁化到處一樣，自由能
\[\mathcal H(r,h,m)=V\big[rm^2+um^4-hm\big].\tag{11.9}\]
〔注〕書上沒有第 9 章的 $\frac12,\frac14$；常數不改變標度（本系列 $m(r,0)=\sqrt{-r/u}$，書上 $\sqrt{-r/2u}$）。\textbf{放大}：改用 $b$ 倍大的單位量長度，體積的數值變成 $Vb^{-d}$。要 $\mathcal H$ 不變、$u$ 也不變，每一項都得自己補回 $b^d$：
$um^4V=u\,m'^4\,Vb^{-d}$ 給 $m'=mb^{d/4}$；$rm^2V=r'm'^2Vb^{-d}$ 給 $r'=rb^{d/2}$；$hmV=h'm'Vb^{-d}$ 給 $h'=hb^{3d/4}$。所以
\[m(r,h)=b^{-d/4}\,m\big(rb^{d/2},hb^{3d/4}\big)\quad\text{對任何 }b\text{ 成立}.\tag{11.11}\]
$b$ 是隨便選的——選 $rb^{d/2}=1$（$b^{-d/4}=r^{1/2}$）就得到\textbf{標度形式}
\[m(r,h)=r^{1/2}f\Big(\frac h{r^{3/2}}\Big),\qquad m(r,0)=\kappa|r|^{1/2},\qquad m(0,h)=\kappa'h^{1/3}.\tag{11.12--11.16}\]
〔注〕(11.13) 是 $r<0$ 的有序相，應寫 $|r|^{1/2}$：$\beta=\frac12$；$r=0$ 時 $\delta=3$。
\textbf{算到底}（$u=1$；方程 $rm+m^3=h$）：$m(-0.25,0)=\sqrt{0.25}=0.5000$；$m(0,0.008)=0.008^{1/3}=0.2000$、$m(0,0.125)=0.5000$。
方程除以 $r^{3/2}$：$y=m/r^{1/2}$、$x=h/r^{3/2}$ 滿足 $y+y^3=x$——$r$ 不見了，這就是 $f$：$f(1)=0.6823$、$f(2)=1$（$1+1=2$）、$f(10)=2$（$2+8=10$）。
\textbf{資料塌縮}（圖 \ref{eleven-fig-collapse}）：$(r,h)=(0.25,0.25)$、$(0.5,0.7071)$、$(1,2)$ 的 $x$ 都是 2（如 $0.25/0.25^{3/2}=2$），所以 $y=1$、$m=r^{1/2}=0.5000,\ 0.7071,\ 1.0000$：三條曲線換軸之後疊成一條。
\begin{errbox}[p.~263，(11.16) 後]
書上說 $r\to0$ 時「$f(h/r^{3/2})\to h^{1/3}r$」。$x$ 大時 $y^3\approx x$，即 $f(x)\approx x^{1/3}$，所以 $f(h/r^{3/2})\to h^{1/3}r^{-1/2}$，乘上 $r^{1/2}$ 才得 $m\to h^{1/3}$。
$r=10^{-4}$、$h=10^{-3}$（$x=1000$）：$f=9.9667$，更正式 10.0000，書上的只有 $10^{-5}$。
\end{errbox}
\textbf{磁化率}（(11.17)）：$\chi=\big(\frac{\partial m}{\partial h}\big)_r\big|_{h=0}=r^{1/2}\cdot r^{-3/2}f'(0)=r^{-1}f'(0)$，所以 $\gamma=1$——只靠「放大後參數怎麼變」就讀出指數。

\textbf{Gaussian 模型}（(11.18)--(11.25)）：現在磁化可以隨位置變，
\[\mathcal H=\int d^dx\Big[\frac r2S(\vec x)^2+\big|\vec\nabla S(\vec x)\big|^2-hS(\vec x)\Big].\tag{11.18}\]
〔注〕梯度項沒有 $\frac12$，和 (11.27) 的 $\frac a2$ 是不同的約定。把座標放大 $\vec x\to b\vec x$：$d^dx\to b^dd^dx$、$\vec\nabla\to b^{-1}\vec\nabla$。要梯度項的樣子不變，$b^{d-2}S^2=S'^2$，即
$S(b\vec x)=b^{(2-d)/2}S'(\vec x)$；代進另外兩項：$rb^dS^2=rb^2S'^2$、$hb^dS=hb^{(d+2)/2}S'$。所以 $r'=b^2r$、$h'=b^{(d+2)/2}h$，
\[m=b^{(2-d)/2}\,m\big(rb^2,hb^{(d+2)/2}\big)\ \xrightarrow{\ b^2=1/r\ }\ m=r^{(d-2)/4}F\Big(\frac h{r^{(d+2)/4}}\Big),\tag{11.19--11.24}\]
$\chi_T=r^{-1}F'$（(11.25)，$\gamma=1$）。〔注〕(11.24) 少了「$m=$」。\textbf{核對}：只有 $q=0$ 的模態感到外場，精確解 $m=h/r$（$F(x)=x$），(11.22)(11.24) 對任何 $b$ 都成立（誤差 $<10^{-12}$）。指數見練習 11.4（$d=4$ 正是平均場的 $\frac12,1,3$）；$d=4$ 時兩種組合完全相同：$\frac{d+2}4=\frac32$、$\frac{d-2}4=\frac12$（$d=3$ 是 $\frac54$、$\frac14$）——第 8 章說過的 $d=4$ 的特殊地位。

\begin{figure}[htbp]
\centering
\ElevenFigCollapse
\caption{Landau 的資料塌縮（$u=1$）。左：$m$ 對 $h$；右：$y=m/r^{1/2}$ 對 $x=h/r^{3/2}$，三條疊成 $y+y^3=x$}\label{eleven-fig-collapse}
\end{figure}

\textbf{揭曉}：只看放大後參數怎麼變就讀出 $\gamma=1$；但這是「被動」的——漲落沒有被積掉。

% =====================================================================
\subsection{實空間重整化：一維 Gaussian 鏈}
\leavevmode\bk{pp.~264--269}\ \anim{W1103}\par\nopagebreak

\textbf{問}：把一半的自旋積掉，剩下的系統還長得一樣嗎？

\textbf{直觀}：\textbf{抽格}（decimation）＝固定偶數格點的自旋，把奇數格點的自旋所有可能的值都積分（加總）掉，只留下它們對偶數格點的「平均影響」。剩下 $N/2$ 個自旋、格距變 2 倍。
如果新的 Hamiltonian 和舊的同一個形式、只是參數換了，就得到參數的\textbf{遞迴關係}：再抽一次、再抽一次……看參數往哪裡跑。\textbf{模型}（連續值的自旋 $S_i$）：
\[\mathcal H[S_i]=\sum_i\Big[\frac r2S_i^2+\frac a2(S_i-S_{i+1})^2\Big].\tag{11.27}\]
$S_3$ 只出現在 $(\frac r2+a)S_3^2-aS_3(S_2+S_4)+\frac a2(S_2^2+S_4^2)$。\textbf{完成平方}（(11.28)--(11.30)，練習 11.5）：用 $\int e^{-AS^2+BS}dS=\sqrt{\pi/A}\,e^{B^2/4A}$，$A=\frac{r+2a}2$、$B=a(S_2+S_4)$，再用 $(S_2+S_4)^2=2(S_2^2+S_4^2)-(S_2-S_4)^2$：
\[\int dS_3\,e^{-(\frac r2+a)S_3^2+aS_3(S_2+S_4)-\frac a2(S_2^2+S_4^2)}=\sqrt{\frac{2\pi}{r+2a}}\exp\Big[-\frac{a^2(S_2-S_4)^2}{2(r+2a)}-\frac{ar(S_2^2+S_4^2)}{2(r+2a)}\Big].\tag{11.30}\]
\textbf{算到底}：$r=0.5$、$a=1$、$S_2=0.3$、$S_4=-0.7$：數值積分 1.224819，閉式 1.224819。
右邊又是「最近鄰耦合＋平方項」：每個偶數格點從左右兩個奇數鄰居各得 $-\frac{ar}{2(r+2a)}S^2$，所以 $\frac{r'}2=\frac r2+\frac{ar}{r+2a}$，
\[r'=\frac{r(r+4a)}{r+2a},\qquad a'=\frac{a^2}{r+2a}\qquad(\text{上例 }r'=0.9,\ a'=0.4).\tag{11.32--11.33}\]
\begin{errbox}[p.~265，(11.31)]
前因子印成 $\big(\frac{2\pi}{r+2a}\big)^{N/2}$：$N/2$ 個奇數格點各給一個 $\sqrt{2\pi/(r+2a)}$，應為 $\big(\frac{2\pi}{r+2a}\big)^{N/4}$。
8 個自旋的 Gaussian 環用行列式精確算 $\ln Z$：照書上多了 $\frac N4\ln\frac{2\pi}{r+2a}=1.8432$（$r=0.5$、$a=1$）。（Ising 的 (11.58) $A^{N/2}$ 是對的，因為 $A$ 不是平方根。）
\end{errbox}

\textbf{重新標度}（(11.34)--(11.38)）：耦合從 $a$ 變成 $a'$，模型「長得不一樣」了。把自旋放大 $S_{2i}\to S_{2i}\sqrt{(r+2a)/a}$（自旋因子），耦合變回 $a$：
\[r''=\frac{r(r+4a)}a,\qquad a''=a,\qquad h''=h\,\frac{r+4a}{\sqrt{a(r+2a)}}\tag{11.36--11.38}\]
（上例 $r''=0.9/0.4=2.25=0.5\times4.5$；外場見練習 11.6）。\textbf{蛛網圖}（Fig.~11.6--11.7，圖 \ref{eleven-fig-chain}）：從 $r_0$ 往上到曲線得 $r_1$，水平走到對角線 $r''=r$（把 $r_1$ 當成新的輸入），再往上……
$r_0=0.05\to0.2025\to0.8510\to4.1282\to33.5553$：越跑越大。\textbf{不動點}是曲線和對角線的交點：$r(r+4a)/a=r$ 給 $r=0$ 與 $r=-3a$。
$r=0$ 處斜率 4 $>1$：一點點偏離就被放大——\textbf{不穩定}的不動點，正是臨界點；〔注〕$r=-3a$ 時 $r+2a<0$，積分發散，不物理。
\begin{errbox}[pp.~267--268，Fig.~11.6--11.7；p.~267 正文]
圖的軸寫 $r/\sqrt a$、$r''/\sqrt a$。$r$ 和 $a$ 同單位，$r''/a=(r/a)(r/a+4)$ 只是 $r/a$ 的函數，軸應為 $r/a$、$r''/a$（書上的曲線就是 $a=1$：$r''=10$ 在 $r=-2+\sqrt{14}=1.7417$）。
正文「the reduced temperature, $t''$, after decimation」的 $t''$ 應為 $r''$。
\end{errbox}

\textbf{小 $r$}（(11.39)--(11.45)，圖 \ref{eleven-fig-chain} 右）：$r\ll a$ 時 $r''\approx4r$、$h''\approx2^{3/2}h$（$h''/h\to2.8284$）、自旋因子 $\approx\sqrt2$，抽 $n$ 次得
$r^{(n)}=4^nr$、$h^{(n)}=2^{3n/2}h$、自旋總因子 $2^{n/2}$。$r=h=10^{-4}$ 時 $k=5$ 次的比只差 1.0086（$h$ 的修正是 $O(r^2)$，比全是 1.0000）。
\textbf{讀出指數}（(11.46)--(11.49)）：舊自旋 $=2^{n/2}\times$ 新自旋，所以 $m=2^{n/2}f(r2^{2n},h2^{3n/2})$ 對每個 $n$ 都成立；選 $2^{2n}=1/r$（$2^{n/2}=r^{-1/4}$）：
\[m=r^{-1/4}F\Big(\frac h{r^{3/4}}\Big),\qquad\Big(\frac{\partial m}{\partial h}\Big)_r=r^{-1}F'\Big(\frac h{r^{3/4}}\Big)\ \Rightarrow\ \gamma=1,\tag{11.48--11.49}\]
正是 (11.24) 取 $d=1$。另一個讀法：長度每次乘 2、$r$ 每次乘 4，所以 $\xi\propto r^{-1/2}$、$\nu=\frac12$。換成 11.6 節的語言：$\Lambda_r=\log_24=2.0000$、$\Lambda_h=\log_22^{3/2}=1.5000$、自旋 $\log_2\sqrt2=0.5000$
（$=\frac{d+2}2$、$\frac{2-d}2$，$d=1$）。
\begin{errbox}[p.~269，(11.46)(11.47)]
書上寫 $m2^{n/2}=f(\cdots)$、$m=2^{-n/2}f(r2^{2n},h2^{3n/2})$。舊自旋是新自旋的 $\sqrt2$ 倍（每次），應為 $m=2^{+n/2}f$。
用精確的 $m=h/r$ 核對（$r=0.01$、$h=0.001$、$n=1$）：照書上 0.0500，更正 0.1000 $=h/r$；(11.48) 的 $r^{-1/4}$ 本來也只和更正版一致。
\end{errbox}

\begin{figure}[htbp]
\centering
\begin{minipage}[c]{0.56\linewidth}\centering\ElevenFigChain\end{minipage}\hfill
\begin{minipage}[c]{0.42\linewidth}\centering\footnotesize\setlength{\tabcolsep}{4pt}\ElevenTabSmallR\end{minipage}
\caption{一維 Gaussian 鏈（$a=1$）。左：Fig.~11.6--11.7（更正的座標 $r/a$）與 $r_0=0.05$ 的蛛網，黑點是不動點 $r=0$（斜率 4）。右：$r=h=10^{-4}$ 抽 $k$ 次，三個比都 $\approx1$}\label{eleven-fig-chain}
\end{figure}

\textbf{揭曉}：積掉一半、重新標度，模型的樣子不變，只有 $r$ 變成 $4r$——由「每次乘 4」讀出 $\nu=\frac12$、$\gamma=1$。

% =====================================================================
\subsection{一維 Ising 的抽格}
\leavevmode\bk{pp.~269--272}\ \anim{W1104}\par\nopagebreak

\textbf{問}：Ising 鏈的奇數自旋加總掉，剩下的自旋之間是什麼耦合？

\textbf{加總一個自旋}（(11.50)--(11.53)）：$Z=\sum_{\{S\}}e^{K\sum_iS_iS_{i+1}}$，$S_i=\pm1$。$S_3$ 只出現在 $e^{KS_3(S_2+S_4)}$：
\[\sum_{S_3=\pm1}e^{K(S_2S_3+S_3S_4)}=2\cosh K(S_2+S_4)=\begin{cases}2\cosh2K,&S_2=S_4\\2,&S_2=-S_4\end{cases}=A\,e^{K'S_2S_4}.\tag{11.52--11.53}\]
\textbf{算到底}（$K=1$）：同號 $2\cosh2=7.5244$、異號 2。要寫成 $Ae^{K'S_2S_4}$：$Ae^{K'}=2\cosh2K$、$Ae^{-K'}=2$，相乘相除（練習 11.7）得
\[A=2\sqrt{\cosh2K},\qquad K'=\frac12\ln\cosh2K\tag{11.54--11.55}\]
$=3.8793$、$0.6625$。新的鏈還是最近鄰 Ising，只是 $K\to K'$。
\begin{errbox}[p.~270，(11.53)]
左邊印成 $\cosh K(S_2+S_4)$，少了 2：$K=1$ 時 $\cosh2=3.7622$，但右邊的 $2\cosh2K=7.5244$。應為 $2\cosh K(S_2+S_4)$。
\end{errbox}
\textbf{遞迴}（(11.62)）：$K_{n+1}=\frac12\ln\cosh2K_n$，$K_1=1$：1.0000、0.6625、0.3501、0.1137、0.0128、0.0002 $\to0$（圖 \ref{eleven-fig-ising}）。
每抽一次，耦合變弱，越看越像高溫：低溫時的長程秩序在更大的尺度上看不見。外場的版本見練習 11.8。
\textbf{配分函數變成連乘}（(11.57)--(11.63)）：每次抽格 $Z_N(K)=A(K)^{N/2}Z_{N/2}(K')$，一直抽下去
\[Z=A(K_1)^{N/2}A(K_2)^{N/4}\cdots A(K_j)^{N/2^j}\cdots .\tag{11.63}\]
\textbf{算到底：4 個自旋的環}（$K=1$）。抽兩次剩 1 個自旋（自己和自己的鍵 $KS_1S_1$ 給 $2e^{K_3}$）：$Z=A(K_1)^2A(K_2)\cdot2e^{K_3}$，$A(1)^2=15.0488$、$A(K_2)=2e^{K_3}=2.8383$，
$Z=15.0488\times2.8383^2=121.2329$。暴力加總 $2^4=16$ 項、轉移矩陣 $(2\cosh K)^4+(2\sinh K)^4=3.0862^4+2.3504^4=90.7140+30.5189$ 都是 121.2329（第 9 章 W902 的 4 自旋環同一個數）；8 個自旋 9160.4387。
\begin{errbox}[p.~271，(11.63)]
最後一個因子印成 $A(K_j)^{N/2^{j+1}}$；$K_1$ 配 $N/2$，應為 $A(K_j)^{N/2^j}$。照書上的指數，4 與 8 個自旋的環得 18.5498、143.2705（應為 121.2329、9160.4387）。
\end{errbox}

\textbf{不動點}（練習 11.9）：令 $x=e^{-2K}$，$e^{2K'}=\cosh2K=\frac{1+x^2}{2x}$ 給 $x'=\frac{2x}{1+x^2}$；$x_0=e^{-2}$ 的蛛網（圖 \ref{eleven-fig-ising} 左）：0.1353、0.2658、0.4965、0.7966、0.9747、0.9997 $\to1$。
不動點 $x=0$（$T=0$，斜率 2，不穩定）與 $x=1$（$T=\infty$，斜率 0，穩定）：除了絕對零度，所有溫度都流向無限高溫。
\textbf{關聯長度}：$x$ 小時 $x'\approx2x$，而抽一次長度單位變 2 倍、$\xi$ 的數值減半：$\xi(2x)=\xi(x)/2$ $\Rightarrow$ $\xi\propto1/x=e^{2J/k_BT}$。
精確 $\xi=-1/\ln\tanh K$，而且 $\tanh K'=\tanh^2K$ 讓 $\xi$ 每次剛好減半；$K=1,2,3$：$\xi=3.6719,\ 27.2960,\ 201.7140$，$e^{2K}/2=3.6945,\ 27.2991,\ 201.7144$。
\textbf{自由能}：$\ln Z/N=\sum_j2^{-j}\ln A(K_j)=\ln(2\cosh K)$；$\ln Z/(NK)=1.1269,\ 1.0091,\ 1.0000$（$K=1,2,5$）$\Rightarrow F=-k_BT\ln Z\to-NJ$（16 個自旋的環用連乘 1.1277）。

\begin{figure}[htbp]
\centering
\begin{minipage}[c]{0.44\linewidth}\centering\ElevenFigXmap\end{minipage}\hfill
\begin{minipage}[c]{0.54\linewidth}\centering\footnotesize\setlength{\tabcolsep}{3.5pt}\ElevenTabIsing\end{minipage}
\caption{一維 Ising。左：$x=e^{-2K}$ 的蛛網（$K=1$ 起）流向 $x=1$。右：$K=1$ 的環，三種算法相同，紅字照書上 (11.63)}\label{eleven-fig-ising}
\end{figure}

\textbf{揭曉}：一維 Ising 的流只有兩個不動點，$T=0$ 不穩定、$T=\infty$ 穩定——所以一維沒有有限溫度的相變。

% =====================================================================
\subsection{動量殼與連續重整化}
\leavevmode\bk{pp.~272--277}\ \anim{W1105}\par\nopagebreak

\textbf{問}：在 Fourier 空間「積掉一半」長什麼樣？

\textbf{直觀}（Fig.~11.8，圖 \ref{eleven-fig-shell} 左）：實空間格距 $c\to2c$ $\Leftrightarrow$ 布里淵區每邊 $Q\to Q/2$。Gaussian 模型在 Fourier 空間各模態互不耦合，所以「積掉外圈殼層」只給一個常數：
拿掉的模態比例 $1-2^{-d}$（$d=2$：$\frac34$、$d=3$：$\frac78$）。
\[S_i=\frac1{\sqrt N}\sum_{\vec q}S(\vec q)e^{i\vec q\cdot\vec R_i},\qquad\mathcal H=\sum_{\vec q}\Big[\frac r2+\alpha(\vec q)\Big]S(\vec q)S(-\vec q)-\sqrt N\,S(0)h.\tag{11.64--11.66}\]
\textbf{算到底}（一維 (11.27)，$a=1$，$N=8$，圖 \ref{eleven-fig-shell} 右的表）：$\frac a2\sum(S_i-S_{i+1})^2=\sum_qa(1-\cos q)|S(q)|^2$，$q=\frac{2\pi k}8$：
$\alpha=0,\ 0.2929,\ 1.0000,\ 1.7071,\ 2.0000$；小 $q$ 近似 $\frac12aq^2=0,\ 0.3084,\ 1.2337,\ 2.7758,\ 4.9348$ 只在 $q$ 小時像。
\begin{errbox}[p.~273，(11.66)(11.67)]
書上 $\alpha(q)=-2a(1-\cos qc)\sim ac^2q^2$。應為 $\alpha(q)=a(1-\cos qc)\sim\frac12ac^2q^2$（$=Jq^2$，$J=\frac12ac^2$）；負號會讓模型不穩定。
隨機 8 個自旋：實空間能量 3.0394＝Fourier；照書上的 $\alpha$ 得 $-4.3508$（負的能量）。
\end{errbox}

\textbf{一次 $b=2$}（(11.68)--(11.73)）：積掉殼層後剩 $|\vec q|<Q/2$；把 $\vec q\to\vec q/2$ 讓布里淵區變回原來大小，$Jq^2\to Jq^2/4$；格點數 $N=2^dN'$；
再令 $S\to2S$ 把 $J$ 變回來：$\frac r2S^2\to\frac r2\cdot4S^2$、$\sqrt NS(0)h=\sqrt{N'}2^{d/2}S(0)h\to\sqrt{N'}2^{(d+2)/2}S(0)h$，
\[r\to2^2r,\qquad h\to2^{(d+2)/2}h,\qquad m(r,h)=b^{(2-d)/2}\,m\big(rb^2,hb^{(d+2)/2}\big)\tag{11.71--11.73}\]
（$d=3$：$h\times5.6569$）——和 (11.22) 一樣，但這次真的積掉了漲落。

\textbf{連續重整化}（(11.74)--(11.83)）：讓 $b=e^{\delta\ell}$ 只比 1 大一點：$\vec q\to\vec q(1-\delta\ell)$、$S\to(1+\delta\ell)S$、$N\to N(1+d\,\delta\ell)$，於是 $r\to r(1+2\delta\ell)$、$h\to h(1+\frac{d+2}2\delta\ell)$：
\[\frac{dr}{d\ell}=2r,\qquad\frac{dh}{d\ell}=\frac{d+2}2h\qquad\Rightarrow\qquad r=r_0e^{2\ell},\quad h=h_0e^{(d+2)\ell/2}.\tag{11.82--11.83}\]
$r_0=h_0=0.01$、$d=3$：$\ell=1$ 得 $(0.0739,\ 0.1218)$，$\ell=2$ 得 $(0.5460,\ 1.4841)$（數值積分相同）；$\ell=\ln2$ 就是一次 $b=2$。
消去 $\ell$：$h=h_0(r/r_0)^{(d+2)/4}\propto|r|^{5/4}$，流線都離開原點（圖 \ref{eleven-fig-shell} 中）——原點是不穩定的不動點，$r$、$h$ 都是相關的方向。
〔注〕(11.78) 的 $N$ 是縮小後的格子以原格子表示。

\textbf{磁化的偏微分方程}（(11.84)--(11.87)）：$S_i=e^{\frac{2-d}2\ell}S_i^{(\ell)}$，所以 $m=e^{\frac{2-d}2\ell}f\big(r(\ell),h(\ell)\big)$ 與 $\ell$ 無關；對 $\ell$ 微分、取 $\ell=0$：
\[0=\Big(-\frac{d-2}2+2r\frac{\partial}{\partial r}+\frac{d+2}2h\frac{\partial}{\partial h}\Big)m(r,h).\tag{11.87}\]
\textbf{代入驗證}：$m=r^{(d-2)/4}F(x)$、$x=h\,r^{-(d+2)/4}$：$2r\partial_rm=\frac{d-2}2m-\frac{d+2}2r^{(d-2)/4}xF'$，$\frac{d+2}2h\partial_hm=\frac{d+2}2r^{(d-2)/4}xF'$，三項一一相消，任何 $F$ 都行
（數值殘差 $<10^{-7}$，$d=1,3,4$）；Landau 的組合 $r^{1/2}f(h/r^{3/2})$ 放進 $d=3$ 就不成立。反過來，(11.87) 的\textbf{特徵線}正是上面的流線：沿著流線 $e^{\frac{2-d}2\ell}m$ 不變，解出來就是 (11.24)。
〔注〕(11.84) 第一行的 $S_i$ 是 $q=0$ 模態對每個格點的貢獻。

\begin{figure}[htbp]
\centering
\begin{minipage}[c]{0.25\linewidth}\centering\ElevenFigShell\end{minipage}\hfill
\begin{minipage}[c]{0.38\linewidth}\centering\ElevenFigFlow\end{minipage}\hfill
\begin{minipage}[c]{0.35\linewidth}\centering\footnotesize\setlength{\tabcolsep}{2.5pt}\ElevenTabAlpha\end{minipage}
\caption{動量殼。左：布里淵區每邊縮一半。中：$d=3$ 的流線（橙點是 $r_0=h_0=0.01$ 的 $\ell=1,2$）。右：$N=8$ 的 $\alpha(q)$ 與書上的值（紅）}\label{eleven-fig-shell}
\end{figure}

\textbf{揭曉}：讓尺度連續地變，重整化變成一組微分方程——參數的「流」。

% =====================================================================
\subsection{不動點與臨界指數}
\leavevmode\bk{pp.~277--284}\ \anim{W1106}\par\nopagebreak

\textbf{問}：參數一直流——會流到哪裡停下來？

\textbf{一般的架構}（(11.88)--(11.92)，Fig.~11.9）：把 $s^d$ 個自旋合成一個\textbf{區塊自旋}，所有參數 $\kappa_i$（$r$、$h$、$u$……）一起變：$\kappa_i\to R_i(\kappa_1,\kappa_2,\ldots)$；連續版
\[\frac{d\kappa_i}{d\ell}=W_i(\kappa),\qquad s=e^\ell,\qquad\frac{d\xi^{(\ell)}}{d\ell}=-\xi^{(\ell)}.\tag{11.89--11.92}\]
長度用 $s$ 倍大的單位量，關聯長度的數值 $\xi^{(\ell)}=\xi e^{-\ell}$ 越來越小。\textbf{臨界超曲面}（$\xi=\infty$ 的參數）被映到自己：$\infty\cdot e^{-\ell}$ 還是 $\infty$。
\textbf{不動點} $W_i(\kappa^*)=0$：流到這裡就停，系統沒有任何長度尺度——$\xi=0$ 或 $\infty$；臨界超曲面上的點都流向一個臨界不動點。

\textbf{線性化}（(11.93)--(11.96)）：在不動點附近令 $\delta\kappa_i=\kappa_i-\kappa_i^*$，$A_{ij}=\partial W_i/\partial\kappa_j$：
\[\frac{d\delta\kappa_i}{d\ell}=\sum_jA_{ij}\delta\kappa_j,\qquad\delta\kappa_i(\ell)=\sum_\alpha C_{i,\alpha}e^{\Lambda_\alpha\ell}.\tag{11.94--11.95}\]
每個特徵方向各自指數地變：$\Lambda>0$ \textbf{相關}（relevant，越流越大，像 $r$ 和 $h$——實驗上要調到 0 才在臨界點）、$\Lambda<0$ \textbf{無關}（irrelevant，越流越小——這些細節被忘掉，不同的模型有相同的指數：\textbf{普適性}）、$\Lambda=0$ \textbf{邊際}（marginal）（註 6）。
\textbf{算到底：示意矩陣}（本系列）$A=\begin{pmatrix}1&2\\2&-2\end{pmatrix}$：特徵多項式 $\lambda^2+\lambda-6=(\lambda-2)(\lambda+3)$。
$A(2,1)^T=(4,2)^T=2(2,1)^T$（相關）、$A(1,-2)^T=(-3,6)^T=-3(1,-2)^T$（無關），所以 $\delta\kappa(\ell)=c_1e^{2\ell}(2,1)+c_2e^{-3\ell}(1,-2)$。
$c_1=0$ 的點留在臨界線（無關方向）上流進不動點；$c_1\ne0$ 的點先靠近（無關的部分消失），再沿 $(2,1)$ 被推開（圖 \ref{eleven-fig-fp} 左）。

\begin{figure}[htbp]
\centering
\begin{minipage}[c]{0.5\linewidth}\centering\ElevenFigToy\end{minipage}\hfill
\begin{minipage}[c]{0.48\linewidth}\centering\footnotesize\setlength{\tabcolsep}{3pt}\ElevenTabExp\end{minipage}
\caption{不動點。左：示意矩陣的流（$\ell\in[0,1]$）。右：(11.105)--(11.108) 與 $\nu=1/\Lambda_r$；後三欄等於第 9 章的指數表}\label{eleven-fig-fp}
\end{figure}

\textbf{$\nu=1/\Lambda_r$}（(11.97)--(11.103)，Fig.~11.10）：溫度稍微偏離臨界，離臨界面的距離 $\propto r$，流 $\ell$ 之後變成 $Kre^{\Lambda_r\ell}$；
流到離臨界面距離 $D$（一個「普通」的距離，那時 $\xi^{(\ell^*)}$ 約是 1 個格距）時停：
\[Kre^{\Lambda_r\ell^*}=D,\quad\xi=e^{\ell^*}\ \Rightarrow\ \xi=\Big(\frac{Kr}D\Big)^{-1/\Lambda_r},\qquad\nu=\frac1{\Lambda_r}.\tag{11.99--11.103}\]
\textbf{算到底}（$K=D=1$、$r=0.01$）：$\Lambda_r=2$：$\ell^*=\frac12\ln100=2.3026$、$\xi=10.0000$；$\Lambda_r=1$：4.6052、100.0000（$\xi=r^{-\nu}$）。
〔注〕(11.98) 的 $C_{i,t}$ 即 $C_{i,r}$，正文的「$re^{y\ell}$」的 $y$ 就是 $\Lambda_r$。

\textbf{四個指數公式}（(11.104)--(11.108)）：每單位體積的自由能，流 $\ell$ 之後體積的數值縮 $e^{\ell d}$：
\[f(r,h)=e^{-\ell d}f\big(re^{\Lambda_r\ell},he^{\Lambda_h\ell}\big).\tag{11.104}\]
選 $e^{\Lambda_r\ell}=1/r$：$f=r^{d/\Lambda_r}\Phi(h\,r^{-\Lambda_h/\Lambda_r})$。對 $r$ 微分兩次得比熱 $\propto r^{d/\Lambda_r-2}$；對 $h$ 微分一次、兩次得 $m\propto r^{(d-\Lambda_h)/\Lambda_r}$、$\chi\propto r^{(d-2\Lambda_h)/\Lambda_r}$；
$r=0$ 時改選 $e^{\Lambda_h\ell}=1/h$：$m\propto h^{d/\Lambda_h-1}$。於是
\[\beta=\frac{d-\Lambda_h}{\Lambda_r},\qquad-\gamma=\frac{d-2\Lambda_h}{\Lambda_r},\qquad2-\alpha=\frac d{\Lambda_r},\qquad\delta=\frac{\Lambda_h}{d-\Lambda_h}\tag{11.105--11.108}\]
（自動滿足 $\alpha+2\beta+\gamma=2$；與第 9 章 Widom 形式的 $\Delta=\Lambda_h/\Lambda_r$ 相同）。
\textbf{Gaussian}（(11.109)--(11.114)）：11.5 節的 $\frac{dr}{d\ell}=2r$、$\frac{dh}{d\ell}=\frac{d+2}2h$ 就是 $\Lambda_r=2$、$\Lambda_h=\frac{d+2}2$，代入得 $\beta=\frac{d-2}4$、$\gamma=1$、$\delta=\frac{d+2}{d-2}$、$\nu=\frac12$——和練習 11.4 相同。
圖 \ref{eleven-fig-fp} 右再放二維 Ising 與二維自迴避（$\Lambda$ 由 $\nu,\beta$ 反推）：正是第 9 章 W903 的指數表；11.3 節的鏈（$\Lambda_r=2$、$\Lambda_h=\frac32$）就是 $d=1$ 的 Gaussian。

\textbf{揭曉}：臨界指數只由不動點附近兩個相關特徵值 $\Lambda_r$、$\Lambda_h$ 決定——下一章把這套用到 O($n$) 模型與自迴避漫步。

% =====================================================================
\subsection{練習 11.1--11.9 解答}
{\small
\begin{enumerate}[label=\textbf{11.\arabic*},itemsep=0.05em]
\item 平面的維度。\quad 答：以面內一點為圓心，$n(R)=\pi R^2$（在邊上 $\frac\pi2R^2$、角上 $\frac\pi4R^2$，前因子不同、指數一樣）$\Rightarrow n(R/3)=\frac19n(R)$ $\Rightarrow3^{-p}=\frac19$、$p=\ln9/\ln3=2$；
複本公式：正方形是 9 個放大 3 倍的小正方形，同樣得 2。
\item Cantor 集。\quad 答：以左端為圓心，$R=\frac13$ 的圓內正好是左邊那一份（右邊那份從 $\frac23$ 開始），而它放大 3 倍就是整個集合 $\Rightarrow n(R/3)=\frac12n(R)$ $\Rightarrow3^{-p}=\frac12$、$p=\ln2/\ln3=0.6309$
（$R=3^{-k}$：1、0.5000、0.2500、0.1250、0.0625）。
\item 複本公式。\quad 答：Koch $\mathcal R=4$、$\mathcal M=3$；Cantor $\mathcal R=2$、$\mathcal M=3$ $\Rightarrow$ (11.6) 給 1.2619、0.6309，和質量－半徑相同；盒子計數 1.2362、0.6254 是第三種估計（有限代數，略偏低）。
\item Gaussian 模型的指數。\quad 答：由 (11.24)：$h=0$、$r<0$ 形式上 $m\propto|r|^{(d-2)/4}$ $\Rightarrow\beta=\frac{d-2}4$；$\chi=r^{-1}F'(0)$ $\Rightarrow\gamma=1$；
$r=0$ 時 $m$ 不能依賴 $r$，要 $F(x)\propto x^{(d-2)/(d+2)}$ $\Rightarrow m\propto h^{(d-2)/(d+2)}$、$\delta=\frac{d+2}{d-2}$。$d=3$：$\frac14,1,5$；$d=4$：$\frac12,1,3$。
〔注〕Gaussian 模型在 $r<0$ 不穩定，$\beta$ 只是標度形式的讀法。
\item 推導 (11.29)--(11.33)。\quad 答：照 11.3 節，$\frac{B^2}{4A}=\frac{a^2(S_2+S_4)^2}{2(r+2a)}$；用 $(S_2+S_4)^2=2(S_2^2+S_4^2)-(S_2-S_4)^2$ 與 $-\frac a2+\frac{a^2}{r+2a}=-\frac{ar}{2(r+2a)}$ 得 (11.30)；
兩邊各一份 $\Rightarrow r'$、$a'$；前因子 $(\frac{2\pi}{r+2a})^{N/4}$（書上 $N/2$）。
\item 加外場。\quad 答：$B=a(S_2+S_4)+h$，$\frac{B^2}{2(r+2a)}$ 的一次項給每個偶數格點兩份 $\frac{ah}{r+2a}$ $\Rightarrow h'=h\frac{r+4a}{r+2a}$，乘自旋因子 $\sqrt{(r+2a)/a}$ 得 (11.38)；
$S^2$ 的係數與 $h$ 無關，(11.36) 不變；常數 $\frac{h^2}{2(r+2a)}$ 只進自由能。核對：$h/r=$ 自旋因子 $\times\,h''/r''$（精確）。
\item 推導 (11.54)(11.55)。\quad 答：$Ae^{K'}=2\cosh2K$、$Ae^{-K'}=2$；相乘 $A^2=4\cosh2K$、相除 $e^{2K'}=\cosh2K$。$K=1$：3.8793、0.6625。
\item 加外場的 Ising。\quad 答：加總 $S_3$ 得 $2\cosh[K(S_2+S_4)+H]$，三種值 $2\cosh(2K\pm H)$、$2\cosh H$ 配 $Ae^{K_1'S_2S_4+K_2'(S_2+S_4)}$：
$e^{4K_2'}=\frac{\cosh(2K+H)}{\cosh(2K-H)}$、$e^{4K_1'}=\frac{\cosh(2K+H)\cosh(2K-H)}{\cosh^2H}$、$A^4=16\cosh(2K+H)\cosh(2K-H)\cosh^2H$；
新外場 $H'=H+2K_2'$（自己的 $H$ 加兩邊各 $K_2'$）。$K=1$、$H=0.1$：3.8897、0.6602、0.0482、0.1964；暴力核對 $Z_8(K,H)=A^4Z_4(K_1',H')$。
〔注〕題目的 $K_1,K_2$ 指 $J/k_BT$、$h/k_BT$，和 (11.62) 的 $K_n$ 撞名。
\item 一維 Ising 的流。\quad 答：(a) $x'=\frac{2x}{1+x^2}$，不動點 $x=0$（不穩定）、$x=1$（穩定）——$\frac12\ln\cosh2K=K$ 只在 $K=0,\infty$。
(b) $x'\approx2x$、$\xi(x')=\xi(x)/2$ $\Rightarrow\xi\propto x^{-1}=e^{2J/k_BT}$。
(c) $K$ 大時 $A(K)\approx\sqrt2e^K$、$K'\approx K-\frac12\ln2$ $\Rightarrow\ln Z/N=\sum_j2^{-j}\ln A(K_j)\to K$、$F\to-NJ$（數值見 11.4 節）。
〔注〕題目 (a) 的「such that that」是重複字；(c) 的「$E=-NJ$」指自由能趨近基態能量 $-NJ$。
\end{enumerate}}

% =====================================================================
\subsection{第 11 章的回顧與筆誤一覽}
\begin{center}
\footnotesize\setlength{\tabcolsep}{4pt}
\begin{tabular}{@{}lll@{}}
\toprule
主題（書上的式子） & 結果 & 動畫\\
\midrule
碎形 (11.1)--(11.6) & Koch $\frac{\ln4}{\ln3}=1.2619$、Cantor 0.6309；普通漫步 2、二維自迴避 $\frac43$ & W1101\\
被動的重整化 (11.7)--(11.25) & $m=r^{1/2}f(h/r^{3/2})$；Gaussian $m=r^{(d-2)/4}F$；$d=4$ 兩者一致 & W1102\\
一維 Gaussian 鏈 (11.26)--(11.49) & $r''=r(r+4a)/a\approx4r$、$h''\approx2^{3/2}h$；$\gamma=1$、$\nu=\frac12$ & W1103\\
一維 Ising (11.50)--(11.63) & $K'=\frac12\ln\cosh2K\to0$、$x'=\frac{2x}{1+x^2}$、$\xi\propto e^{2K}$ & W1104\\
動量殼 (11.64)--(11.87) & $\frac{dr}{d\ell}=2r$、$\frac{dh}{d\ell}=\frac{d+2}2h$；(11.24) 滿足 (11.87) & W1105\\
不動點 (11.88)--(11.114) & $\nu=1/\Lambda_r$、$\beta=(d-\Lambda_h)/\Lambda_r$；Gaussian $\Lambda_r=2$、$\Lambda_h=\frac{d+2}2$ & W1106\\
\bottomrule
\end{tabular}
\end{center}
\textbf{筆誤一覽}（給 2004 年版；〔注〕見各節）：下表在文中都用紅框標出，每一條都用程式核對過。
\begin{center}
\footnotesize\renewcommand{\arraystretch}{0.9}
\ElevenTabErrata
\end{center}
